\subsection{Solucion desarrollada}

\begin{respuesta}{(a) Limites para promedio y recorrido}
Calculamos media y recorrido de cada submuestra:
\[
\bar X=(0.45,\ 0.50,\ 0.60,\ 0.50),
\qquad
R=(0.30,\ 0.20,\ 0.20,\ 0.00).
\]
Por lo tanto:
\[
\bar{\bar X}=\frac{0.45+0.50+0.60+0.50}{4}=0.5125,
\qquad
\bar R=\frac{0.30+0.20+0.20+0}{4}=0.175.
\]
Reemplazando:
\[
LCS_{\bar X}=0.5125+0.729(0.175)=0.6401,
\qquad
LCI_{\bar X}=0.5125-0.729(0.175)=0.3849.
\]
Para la carta de recorridos:
\[
LCS_R=2.282(0.175)=0.3994,\qquad LCI_R=0(0.175)=0.
\]
Por lo tanto, los limites para el promedio son:
\[
\boxed{LCI_{\bar X}=0.3849},\qquad
\boxed{LC_{\bar X}=0.5125},\qquad
\boxed{LCS_{\bar X}=0.6401}.
\]
Los limites para la diferencia entre maximo y minimo, es decir, para el recorrido \(R\), son:
\[
\boxed{LCI_R=0},\qquad
\boxed{LC_R=0.175},\qquad
\boxed{LCS_R=0.3994}.
\]
\end{respuesta}

\begin{respuesta}{(b) Diagnostico con la nueva muestra}
La muestra nueva tiene:
\[
\bar X_5=\frac{0.4+0.7+0.5+0.9}{4}=0.625,
\qquad
R_5=0.9-0.4=0.5.
\]
El promedio esta dentro del intervalo \([0.3849,0.6401]\), pero el recorrido no:
\[
0.5>0.3994.
\]
Por lo tanto, el proceso no esta bajo control. La causa observable es exceso de dispersion dentro de la muestra, no un desplazamiento claro de la media.
\end{respuesta}

\begin{respuesta}{(c) Plan de muestreo para controlar insumos}
Se usa un plan de muestreo simple \((n,c)\): se inspeccionan \(n\) unidades y se acepta el lote si aparecen como maximo \(c\) defectuosas. Los parametros del enunciado se interpretan asi:
\[
\begin{array}{c|c|c}
\text{Parametro} & \text{Valor} & \text{Interpretacion} \\
\hline
AQL & 0.02 & \text{calidad aceptable: }2\%\text{ defectuoso} \\
LPTD & 0.08 & \text{calidad rechazable: }8\%\text{ defectuoso} \\
\alpha & 0.05 & \text{riesgo de rechazar un lote bueno} \\
\beta & 0.10 & \text{riesgo de aceptar un lote malo}
\end{array}
\]

Primero se calcula la razon que se busca en la tabla:
\[
\frac{LPTD}{AQL}
=\frac{0.08}{0.02}
=4.
\]
Se usa la tabla correspondiente a \(\alpha=0.05\) y \(\beta=0.10\). En esa tabla se busca el valor de \(\frac{LPTD}{AQL}\) cercano a \(4\):
\[
\begin{array}{c|c|c}
c & LPTD/AQL & n\cdot AQL \\
\hline
3 & 4.89 & 1.366 \\
4 & 4.057 & 1.97 \\
5 & 3.549 & 2.613
\end{array}
\]
Como \(4.057\) es el valor mas cercano a \(4\), se toma:
\[
c=4,
\qquad
n\cdot AQL=1.97.
\]
Despejando \(n\):
\[
n=\frac{1.97}{0.02}=98.5\approx 99.
\]
Por lo tanto, el plan de muestreo es:
\[
\boxed{n=99,\qquad c=4.}
\]
La regla final es: tomar una muestra de \(99\) unidades; si hay \(4\) defectuosas o menos, se acepta el lote; si hay \(5\) defectuosas o mas, se rechaza.
\end{respuesta}
